% --- VII. Conclusion ---
\section{Conclusion}
\label{sec:conclusion}
In this project, our team built \ossim, a kernel simulator in which CPU scheduling, synchronization, paging and disk scheduling run together. Starting from the scaffold, we added four CPU schedulers, deadlock detection, two page-replacement policies, three disk schedulers and a workload generator. We checked every policy against examples solved by hand, and all 45 automatic tests pass.

The experiments show two classic trade-offs. For CPU scheduling, SJF and SRTF give the lowest waiting time, while RR and MLFQ give the lowest response time, at every level of I/O. For memory, CPU utilization drops sharply once the frames can no longer hold the pages that the processes are using, and the choice of replacement policy matters much less than the number of processes in memory. \todo{One sentence with the main finding of the theme experiment.}

The simulator also has some limitations, which are possible future work:
\begin{itemize}
    \item SJF and SRTF know the exact length of the next burst. A real OS must predict it, for example with exponential averaging~\cite[Sec.~5.3.2]{silberschatz}.
    \item PRIO has no \emph{aging}, so a process with a low priority can starve.
    \item A page that has just been loaded is not protected until its owner uses it. With very little memory, this can make the system almost stop, as Experiment~2 shows.
\end{itemize}
